\documentclass[11pt]{article}
\usepackage[T1]{fontenc}
\usepackage{lmodern}
\usepackage[letterpaper,margin=1in]{geometry}
\usepackage{amsmath,amssymb,amsthm,mathtools,microtype}
\usepackage[hidelinks]{hyperref}
\hypersetup{pdftitle={A nonzero scalar shift charge in perturbative de Sitter gravity}}
\newtheorem{theorem}{Theorem}
\newtheorem{lemma}[theorem]{Lemma}
\newcommand{\HH}{\mathcal H}
\newcommand{\DD}{\mathcal D}
\newcommand{\FF}{\mathcal F}
\newcommand{\Tr}{\operatorname{Tr}}
\newcommand{\Sym}{\operatorname{Sym}}
\newcommand{\SO}{\mathrm{SO}}
\newcommand{\dd}{\,\mathrm d}
\newcommand{\vac}{\Omega_{\rm vac}}
\numberwithin{equation}{section}
\title{A nonzero scalar shift charge in perturbative de Sitter gravity}
\author{}
\date{}
\begin{document}
\maketitle
\begin{abstract}
We retain the homogeneous Schr\"odinger mode of a massless minimally
coupled scalar when imposing the quadratic linearization-stability
constraints of Einstein--scalar gravity on global \(\mathrm{dS}_4\).
The connected de Sitter group acts on the full scalar Weyl algebra:
boosts displace the radiation by an amount proportional to the shift
charge, and antipodal symmetry makes the displacement cocycle isotropic.
On each compact charge interval away from zero, a sector with eight
physical transverse-traceless gravitons has a dense invariant seed
domain on which Haar averaging is absolutely convergent and positive
semidefinite. We prove positivity by splitting the graviton tensor
factors into two groups of four and writing the averaged form as a
trace pairing of positive matrix convolution operators. An explicit
normalizable seed has strictly positive averaged norm and physical
shift-charge spectrum supported away from zero. The construction
imposes all ten connected de Sitter generators at quadratic
constraint order, without using a sharp charge eigenstate.
\end{abstract}

\section{Setting and result}

Consider Einstein gravity with cosmological constant \(3/R^2>0\)
and a real minimally coupled scalar with no potential, on the exact
constant-scalar global de Sitter background
\[
 \bar g=R^2(-\mathrm dt^2+\cosh^2t\,\mathrm d\Omega_3^2),
 \qquad \bar\phi=0 .
\]
The time \(t\) is dimensionless; the scalar kinetic term is
\(-\tfrac12\int |d\phi|_g^2\dd{\rm vol}_g\). We set \(\hbar=1\).
Let \(G=\SO_0(4,1)\) act on solutions by inverse pullback and
\(K=\SO(4)\) be its maximal compact subgroup. After ordinary local
gauge reduction, the linear phase space is the direct sum of the
scalar solutions and physical spin-two solutions. Compact Cauchy
surfaces also enforce ten integrated quadratic constraints, one for
each Killing generator. They require the simultaneous scalar and
graviton action, not either action separately
\cite{higuchi,higuchiii,mmfree,my}.

The group-averaging form is
\begin{equation}
 {\cal G}(\psi,\chi)=\int_G
             \langle\psi,U(g)\chi\rangle\,\dd g .
 \label{eq:ga}
\end{equation}
The inner product is conjugate linear in its first argument and
\(\dd g\) is a fixed positive Haar normalization. The averaged objects
are functionals on a test domain, not normalizable invariant vectors
in the kinematic Hilbert space. Absolute convergence alone does not
imply positivity. Marolf and Yu establish positivity for a broad class
of ordinary de Sitter Fock states, explicitly setting aside the
non-Fock homogeneous massless scalar mode in dimension greater than
two \cite[Sec.~6]{my}. A positive averaging result retaining the
scalar zero mode is known in \(\mathrm{dS}_2\) \cite{mm2}.
The coupling between the scalar zero mode and radiation under de
Sitter boosts is itself known, including in the compact axion
quantization \cite[Secs.~2.3--2.6]{axion} and in the noncompact
mode analysis of \cite[Sec.~II, Eqs.~(31)--(40)]{gh}. Here the new question is the positive
\emph{gravitational} constraint projection at nonzero charge.

\begin{theorem}[A charged positive sector]\label{thm:main}
Let \(J=[p_-,p_+]\) with \(0<p_-<p_+<\infty\).
The connected de Sitter action on the full scalar Weyl algebra is
unitarily implemented in the regular representation
\(L^2(\mathbb R,\dd p;\FF_{\rm rad})\), where \(p\) is conserved.
Let \(\HH_{\rm TT}\) be the positive one-particle Hilbert space
of the physical graviton and set
\[
 \HH_{J,8}
 =L^2(J,\dd p;\FF_{\rm rad})\otimes \Sym^8\HH_{\rm TT}.
\]
There is a dense invariant linear domain \(\Phi\subset\HH_{J,8}\)
on which \eqref{eq:ga} converges absolutely for every pair and is
positive semidefinite. Its positive quotient and completion carry a
self-adjoint shift charge with spectral values in \(J\). An explicit
normalizable seed has strictly positive averaged norm and charge law
supported in \(J^\circ\). Its averaged functionals satisfy the ten
connected quantum quadratic constraints on a dense smooth test core.
\end{theorem}

Eight gravitons are sufficient, not claimed minimal.
The fixed-\(J\) sector is \(G\)-invariant, but does not carry
the translations conjugate to \(p\); the full scalar Weyl
representation is constructed before this restriction.

\section{The full scalar Weyl representation}
\subsection{The homogeneous mode and the radiative quotient}

Let \(v_3\) be the volume of the unit \(S^3\). For smooth real scalar
solutions define the conserved forms
\begin{align}
 p(u)&=R^2\cosh^3t\int_{S^3}\partial_tu\,\dd\Omega,\\
 \Omega(u,w)&=R^2\cosh^3t
       \int_{S^3}(u\partial_tw-w\partial_tu)\,\dd\Omega .
 \label{eq:symp}
\end{align}
For \(e=1\) and \(m(t)=(R^2v_3)^{-1}\int_0^t\cosh^{-3}s\,\dd s\),
we have \(p(m)=\Omega(e,m)=1\). Each solution is uniquely
\(u=q e+p m+v\), where \(q\) is its mean at \(t=0\), \(p=p(u)\),
and \(v\) has both initial means zero. Let \(V\) be the corresponding
radiative solution space. It is symplectically orthogonal to
\(\operatorname{span}(e,m)\) and carries \(\Omega_V\).
The homogeneous sector is a canonical plane, not an oscillator.
The formal invariant \(p=0\) vacuum of the standard zero-mode
discussion \cite[Sec.~IV]{kg} is not a normalizable seed here.

The positive one-particle structure on \(\ker p/\mathbb R e\cong V\)
comes from the nonconstant Euclidean modes. For instance, in
conformal time \(\tan\eta=\sinh t\), degree \(j\ge1\), and \(k=j+1\),
an orthonormal real sphere harmonic has positive mode
\[
 u_j(\eta)=
 \frac{(k\cos\eta+i\sin\eta)e^{-ik\eta}}
 {\sqrt{2R^2 k(k^2-1)}} .
\]
Indeed the scalar mode equation is
\(u_j''+2\tan\eta\,u_j'+j(j+2)u_j=0\); at the neck the value is
\(k/\sqrt{2R^2k(k^2-1)}\) and the first derivative is
\(-i(k^2-1)/\sqrt{2R^2k(k^2-1)}\), giving unit positive
Klein--Gordon norm. Continuation to \(\eta=-is\) for
\(s\to+\infty\) decays as \(e^{-js}\), the regular Euclidean mode.
Its real solution norm at the neck is equivalent to the sum of
\(k\) times squared configuration amplitudes and \(k^{-1}\) times
squared velocity amplitudes; smooth \(V\)-data have finite norm.

The massless limit of the scalar boost recurrence
\cite[Eqs.~(16)--(21), Sec.~2.2]{mmfree} preserves the positive
nonconstant modes modulo \(\mathbb R e\). Specifically, at fixed
\(j\ge1\) the normalization factors
\(\Gamma(j-\sigma)\Gamma(j+\sigma+3)\) are finite at \(\sigma=0\).
The only possible transition to the homogeneous mode is \(j=1\)
to \(j=0\). Its coefficient is \(O(\sqrt{-\sigma})\), and the
divergent leading piece of the normalized homogeneous mode is
\(O((-\sigma)^{-1/2})\) times the constant solution; the product
is zero in the quotient. The remaining boost coefficients
preserve positive frequency and are \(O(j+1)\).
On finite-mode vectors, the \(n\)-fold boost derivative has bound
\(C^n(j+n)!/j!\), with a polynomial factor if any fixed number
of additional Sobolev weights is used. This proves local
convergence of the series in the one-particle norm and the smooth
solution norms. Uniqueness of evolution by the boost differential
equation identifies the series with the solution pullback modulo
constants. Iteration and rotations give the full action \(T_g\),
unitary by the conserved symplectic form and positive structure.
Let \(\|\cdot\|_\mu\) be its real norm in the convention in which
the radiation vacuum Weyl amplitude is
\(\exp(-\|v\|_\mu^2/4)\).

\subsection{An isotropic boost cocycle}

Because \(p(u)=\Omega(e,u)\) and \(e\) is fixed, \(p\) is invariant.
The de Sitter antipode \(\mathcal A\) in embedding coordinates
commutes with \(G\) and is anti-symplectic. It preserves \(V\).
The solution \(m\) is antipodally odd, so \(gm-m\) is odd, has
charge zero and has zero mean at the neck. Thus
\begin{equation}
 gm=m+c_g,\qquad gv=T_gv+\ell_g(v)e,\qquad c_g\in V.
 \label{eq:split}
\end{equation}
The odd subspace of an anti-symplectic involution is isotropic.
Since the induced antipode commutes with \(T_g\), the group law
and invariance of \(\Omega\) give
\begin{equation}
 c_{gh}=c_g+T_gc_h,\qquad
 \Omega_V(c_g,T_gc_h)=0,\qquad
 \ell_g(v)=\Omega_V(c_g,T_gv).
 \label{eq:cocycle}
\end{equation}
The last identity is \(\Omega(gm,gv)=0\). Each \(c_g\) is smooth
and depends continuously on finite \(g\) in the radiation norm.
The cocycle is nonzero: an axial boost of \(m\) differentiated
at the neck gives the \(j=1\) spatial configuration profile
\(\pm(R^2v_3)^{-1}\cos\vartheta\), where \(\vartheta\) is the
polar angle of the boost axis. This follows directly from the
normal component of the boost at the neck and \(m'(0)\).

\subsection{Weyl covariance including the charge translations}

Use \(W_V(v)W_V(w)=e^{-i\Omega_V(v,w)/2}W_V(v+w)\).
In \(\HH_\phi=L^2(\mathbb R,\dd p;\FF_{\rm rad})\), let
\(\widehat p=p\) and \(\widehat q=i\partial_p\) on the standard
domains. The full Weyl representation is
\begin{equation}
 (\mathcal W(a,b,v)\Psi)(p)
       =e^{ia(p+b/2)}W_V(v)\Psi(p+b).
 \label{eq:fullWeyl}
\end{equation}
This implements the symplectic form
\(a b'-b a'+\Omega_V(v,v')\).
Define the strongly continuous unitary action
\begin{equation}
 (U_\phi(g)\Psi)(p)=W_V(-pc_g)\Gamma(T_g)\Psi(p).
 \label{eq:u}
\end{equation}
The Weyl multiplier in the composition vanishes by
\eqref{eq:cocycle}, so this is not merely a projective action.
Strong continuity follows on each fiber, then by dominated
convergence in \(p\). Conjugating \eqref{eq:fullWeyl} gives
\[
 U_\phi(g)\mathcal W(a,b,v)U_\phi(g)^*
   =\mathcal W(a+\ell_g(v),b,T_gv+b c_g).
\]
Indeed \(W_V(-pc_g)W_V(T_gv)W_V((p+b)c_g)\) has phase
\(e^{i(p+b/2)\ell_g(v)}\). This proves covariance for the
original full scalar algebra, not just an externally prescribed
charge fiber. A scalar radiation-vacuum charge packet \(f\) has
\begin{equation}
 \langle f\vac,U_\phi(g)f\vac\rangle
    =\int_{\mathbb R}|f(p)|^2
      e^{-p^2\|c_g\|_\mu^2/4}\dd p .
 \label{eq:scalar-vac}
\end{equation}
For a nonzero packet it is strictly positive for all finite
\(g\); the charged scalar vector itself is not invariant.

\section{The averaging domain}

We use the usual positive two-helicity graviton modes and their
action on the physical temporal transverse-traceless quotient
\cite[Secs.~3.3--3.4]{mmfree}. Let \(\DD_{\rm TT}\) be the
space of finite sums of the spherical harmonic multiplets.
In the Cartan decomposition \(g=k_1a_\lambda k_2\),
\(k_i\in K\), the finite-mode matrix elements obey
\begin{equation}
 |\langle v,U_{\rm TT}(g)w\rangle|
             \le C(v,w)e^{-\lambda},\qquad v,w\in\DD_{\rm TT}.
 \label{eq:bound}
\end{equation}
For a fixed pair of harmonics this follows from the explicit
boost coefficients of \cite[Eqs.~(36)--(42), (125)]{mmfree}.
Their two large-rapidity exponents are \(\sigma-k\) and
\(-3-\sigma-k\). The tensor list has \(\sigma\in\{-2,-1,0\}\)
and \(k\ge2\), bounding both exponents by \(-2\), whether
or not a particular listed family occurs in four dimensions.
The weaker \eqref{eq:bound} absorbs fixed polynomial prefactors
and applies at all \(\lambda\ge0\). Rotations preserve the finite multiplets,
so the constant can be taken uniform in both rotations.
The Cartan Haar density for \(\SO_0(4,1)\) is proportional to
\((\sinh\lambda)^3\dd\lambda\).

Let
\(\Phi_0=L^2(J,\dd p;\FF_{\rm rad})\odot\Sym^8\DD_{\rm TT}\),
where \(\odot\) denotes algebraic tensor span. Embed symmetric
powers isometrically into ordinary tensor powers.
With \(U=U_\phi\otimes U_{\rm TT}^{\otimes8}\) restricted to the
symmetric subspace, define
\(\Phi=\operatorname{span}\{U(g)\psi:\psi\in\Phi_0,\ g\in G\}\).
It is dense and invariant. For \(\psi,\chi\in\Phi_0\), finite tensor
expansion and \eqref{eq:bound} give
\[
 |\langle\psi,U(g)\chi\rangle|\le C_{\psi,\chi}
                                     e^{-8\lambda(g)}.
\]
The scalar contributes only the product of its two norms, regardless
of how many radiation quanta a boost produces.
The same estimate holds on each charge fiber before integrating
the products of the scalar fiber norms. Translated seeds preserve
the exponential rate, with changed constants, by
\(\lambda(h^{-1}gk)\ge\lambda(g)-\lambda(h)-\lambda(k)\)
on \(G/K\). This proves absolute convergence on \(\Phi\)
and permits integration over charge by Fubini. Unimodular Haar
invariance also gives
\({\cal G}(U(h)\psi,U(k)\chi)={\cal G}(\psi,\chi)\)
independently for each argument.

\section{Positivity from two convolution operators}

Consider a unimodular group with fixed Haar measure. Let
\(R_g\xi(x)=\xi(xg)\) be the right regular action. For a finite
family \(x_i\) in a unitary representation \(L\), write
\(F_{ij}(g)=\langle x_i,L(g)x_j\rangle\).
If the entries are in \(L^1(G)\), define the bounded operator
\(R(F)=\int F(g)\otimes R_g\dd g\). Its kernel at \(x,y\)
has entries \(F_{ij}(x^{-1}y)\), which are the Gram kernel of
the vectors \(L(x)x_i\). It follows that \(R(F)\ge0\), first
on compactly supported vector tests and then everywhere.
Entrywise conjugation preserves positive type.

For the right group von Neumann algebra use the canonical
Plancherel trace \(\tau\). The unimodular group Hilbert-algebra
identity is
\[
 \tau(R(f)^*R(g))=\int_G\overline{f(x)}g(x)\dd x,
                       \qquad f,g\in L^1\cap L^2.
\]
Inversion interchanges the right and left conventions
\cite{plancherel}. The right involution is
\(f^*(g)=\overline{f(g^{-1})}\). In finite matrix entries we use
\(\Tr_N\otimes\tau\).

\begin{lemma}[Tensor product positivity on a decaying core]
\label{lem:convolution}
Let \(\psi=\sum_{i=1}^{N}x_i\otimes y_i\) and let \(L,M\) be
unitary representations. Suppose both
\(F_{ij}(g)=\langle x_i,L(g)x_j\rangle\) and
\(H_{ij}(g)=\langle y_i,M(g)y_j\rangle\) belong entrywise to
\(L^1(G)\cap L^2(G)\). Then
\(\int_G\langle\psi,(L(g)\otimes M(g))\psi\rangle\dd g\)
converges absolutely and is nonnegative.
\end{lemma}

\begin{proof}
Absolute convergence follows entrywise by Cauchy--Schwarz.
The positive operators \(A=R(F)\), \(B=R(\overline H)\) are in
trace-\(L^2\) for the matrix Plancherel trace. Their square roots
are trace-\(L^4\), whence \(B^{1/2}A^{1/2}\) is trace-\(L^2\)
and the trace of \(AB\) is the squared trace-\(L^2\) norm of
that operator. In detail, the Plancherel identity and
\(H_{ji}(g^{-1})=\overline{H_{ij}(g)}\) give
\begin{align*}
 (\Tr_N\otimes\tau)(AB)
   &=\sum_{i,j}\int_G F_{ij}(g)
                              \overline{H_{ji}(g^{-1})}\dd g\\
   &=\sum_{i,j}\int_G F_{ij}(g)H_{ij}(g)\dd g\\
   &=\int_G\langle\psi,(L(g)\otimes M(g))\psi\rangle\dd g .
\end{align*}
No convolution distribution is evaluated at one point.
\end{proof}

Group the eight graviton factors of the ambient tensor power
into four plus four, putting the scalar space in the first group.
This is an isometric estimate on the actual symmetric tensors,
not a new particle labelling in the physical model.
Every seed in \(\Phi_0\) is a finite sum across these two groups.
The two coefficient arrays decay as \(e^{-4\lambda}\) and hence
lie in \(L^1\cap L^2\), including when the scalar seed is entangled
with graviton factors. Lemma~\ref{lem:convolution} proves positivity
on \(\Phi_0\). Independent translate invariance of \eqref{eq:ga}
reduces the positive form on a finite sum of translates to the
form on the unshifted sum. Thus \(\cal G\) is positive
semidefinite throughout \(\Phi\).

\section{The physical charge and the constraints}

\subsection{A strictly positive charged seed}

On positive-frequency gravitons set
\(\mathcal J h=\overline{\mathcal A^*h}\).
The antipode reverses time orientation, so both factors in
\(\mathcal J\) reverse the sign of the complex Klein--Gordon
form. Their product is an antiunitary involution.
It preserves the physical positive modes: for a real tensor
harmonic the Euclidean normalized time mode has a real value
and purely imaginary normal derivative at the neck, up to a
common phase, by reflection and conjugation of the regular
mode of \cite[Eq.~(121)]{mmfree}. The spatial antipode is a
real orthogonal involution on the tensor multiplet.
Pure gauge fields map to pure gauge fields; \(\mathcal J\)
commutes with \(G\) on the quotient. Choose a unit real
antipodal harmonic with the Euclidean positive time mode and
phase it by \(1\) or \(i\) to get \(v\in\DD_{\rm TT}\) with
\(\mathcal Jv=v\).
Thus \(a(g)=\langle v,U_{\rm TT}(g)v\rangle\) is real and
\(a(1)=1\).

For a unit \(f\in C_c^\infty(J^\circ)\), the normalizable seed
\(\psi=f\vac\otimes v^{\otimes8}\) gives
\begin{equation}
 {\cal G}(\psi,\psi)=\int_J |f(p)|^2\rho(p)\dd p,\qquad
 \rho(p)=\int_G a(g)^8
                    e^{-p^2\|c_g\|_\mu^2/4}\dd g .
 \label{eq:density}
\end{equation}
The integrand is nonnegative and strictly positive near the identity.
The bound \eqref{eq:bound} and dominated convergence give a
finite continuous \(\rho\). Restricting the integral to a compactly
contained identity neighborhood also proves a strictly positive
lower bound uniform on \(J\).
The nonzero state uses neither a formal \(\delta(p-p_0)\) nor a
presumed invariant charged scalar vacuum.

\subsection{The descended spectral charge}

Multiplication by \(b\in L^\infty(J)\) preserves \(\Phi\),
commutes with \(U\), and has \({\cal G}\)-adjoint multiplication
by \(\bar b\). For \(b\) the indicator of a Borel set, it and
its complement give two orthogonal nonnegative parts of the
quadratic form. They descend to bounded orthogonal projections
on the positive quotient and completion. Countable additivity
holds strongly on the core by dominated convergence in the joint
charge and Haar integrals, using the fiber version of the
matrix-element bound of Sec.~3. It extends to the completion by
the projection norm bounds. This defines the PVM of the bounded
self-adjoint charge on this sector. The vector of \eqref{eq:density}
has normalized charge distribution
\[
 \frac{|f(p)|^2\rho(p)\dd p}
      {\int_J|f(p')|^2\rho(p')\dd p'} .
\]
This is supported away from zero. The conjugate \(\widehat q\)
need not be an invariant observable on the quotient.

\subsection{All ten integrated constraints}

On a compact slice the constraint Hamiltonian \(C_X\) smeared
with the lapse and shift of a background Killing field \(X\) has
a vanishing Hamiltonian vector at the background and no
boundary generator. In a local canonical chart,
differentiate \(dC_X(w)=\Omega(w,X_{C_X})\) at the background:
its Hessian on the linear physical solutions is
\(\Omega(w,B_Xu)\), where \(B_X\) is the infinitesimal Killing
action, up to the common sign for inverse pullback.
The integrated quadratic Taub constraint is consequently the
moment map \(\tfrac12\Omega(u,B_X u)\) for the linear
symplectic action, summed over scalar and graviton
\cite{higuchi,mmfree,my}. This prescription does not omit the
zero-mode part of the scalar generator.
Explicitly, if \(C_X^{\rm rad}=\partial_s c_{\exp(sX)}|_{s=0}\)
and \(B_X^{\rm rad}\) generates \(T_{\exp(sX)}\), then the scalar
moment map on \(u=qe+pm+v\) is
\[
 \tfrac12\Omega_V(v,B_X^{\rm rad}v)
                          -p\,\Omega_V(C_X^{\rm rad},v).
\]
Indeed the constant component of \(B_Xu\) is
\(\Omega_V(C_X^{\rm rad},v)e\) and its \(V\)-component is
\(p C_X^{\rm rad}+B_X^{\rm rad}v\).
The mixed term is nontrivial for boosts; the Weyl displacement in
\eqref{eq:u} implements it. The unitary representation has no
extra central term.

Let the smooth core consist of finite \(G\)-translates of finite sums
with smooth compact charge profiles, finite-particle radiation in
finite harmonic multiplets, and finite graviton multiplets.
It is dense in \(\HH_{J,8}\).
Its Lie derivatives are again in this class: one radiation or
graviton generator moves finitely many angular degrees, and the
derivative of the displacement uses the \(j=1\) finite radiation
mode and bounded \(p\) on each seed, or vanishes for rotations.
On a translated seed use the adjoint action on the generator.
The absolute convergence bound applies also with such a derivative
and with a short bounded group translation of either argument.
Differentiating Haar invariance under the integral therefore gives
\[
 {\cal G}(\psi,dU(X)\chi)=0
\]
on this core for every \(X\in\mathfrak{so}(4,1)\), and likewise in
the first slot. This is the weak constraint equation on the averaged
functionals. It neither presumes that the invariant vectors are
normalizable in \(\HH_{J,8}\) nor that the nonlinear constraints
have been constructed. This completes the proof of
Theorem~\ref{thm:main}.

\section{Scope}

The result is a positive constraint quotient for a specified
non-Fock matter sector in perturbative gravity. The tensor split
proves positivity; it is not a factorization of spatial regions.
The charge survives all ten connected linearization-stability
constraints at quadratic order. The construction neither selects
a nonperturbative vacuum, nor determines the minimal graviton
number, a general chart of other charge sectors, or higher-order
constraint dynamics.

\begin{thebibliography}{9}
\bibitem{higuchi}
A.~Higuchi, Quantum linearization instabilities of de Sitter
space-time. I, \emph{Class. Quantum Grav.}\ \textbf{8},
1961--1981 (1991).

\bibitem{higuchiii}
A.~Higuchi, Quantum linearization instabilities of de Sitter
space-time. II, \emph{Class. Quantum Grav.}\ \textbf{8},
1983--2004 (1991),
\href{https://doi.org/10.1088/0264-9381/8/11/010}{doi:10.1088/0264-9381/8/11/010}.

\bibitem{mmfree}
D.~Marolf and I.~A. Morrison, Group averaging for de Sitter free
fields, \emph{Class. Quantum Grav.}\ \textbf{26}, 235003 (2009),
\href{https://arxiv.org/abs/0810.5163}{arXiv:0810.5163}.

\bibitem{my}
D.~Marolf and X.~Yu, On the positivity of de Sitter group
averaging (2026), \href{https://arxiv.org/abs/2610.24073}{arXiv:2610.24073}.

\bibitem{mm2}
D.~Marolf and I.~A. Morrison, Group averaging of massless scalar
fields in \(1+1\) de Sitter (2008),
\href{https://arxiv.org/abs/0808.2174}{arXiv:0808.2174}.

\bibitem{kg}
K.~Kirsten and J.~Garriga, Massless minimally coupled fields in
de Sitter space: \(O(4)\)-symmetric states versus de Sitter invariant
vacuum, \emph{Phys. Rev. D}\ \textbf{48}, 567--577 (1993),
\href{https://arxiv.org/abs/gr-qc/9305013}{arXiv:gr-qc/9305013}.

\bibitem{gh}
J.~Gibbons and A.~Higuchi, Removing the Faddeev--Popov zero modes
from Yang--Mills theory in spacetimes with compact spatial sections,
\emph{Phys. Rev. D}\ \textbf{91}, 024006 (2015),
\href{https://arxiv.org/abs/1410.7830}{arXiv:1410.7830}.

\bibitem{axion}
V.~A. Letsios and S.~Vitouladitis, Axions on de Sitter space
(2026), \href{https://arxiv.org/abs/2606.28858}{arXiv:2606.28858}.

\bibitem{plancherel}
J.~Dixmier, \emph{Von Neumann Algebras}, North-Holland (1981),
sections on the Hilbert algebra of a locally compact unimodular group.
\end{thebibliography}
\end{document}
